\section{Introduction}
\label{sec:intro}

Deep neural networks are known to be vulnerable to \emph{backdoor attacks}: an
adversary poisons a fraction of the training data so that the learned model
behaves normally on clean inputs but misclassifies any input bearing a trigger
to a target label~\cite{gu2017badnets}. Most early attacks
\emph{flip the label} of poisoned samples, which is easy to spot under visual
inspection. \textbf{Clean-label} attacks~\cite{turner2018labelconsist,chen2019blended}
are stealthier: poisoned samples keep their \emph{true} (target-class) label, so
the attack is invisible to label-level auditing. Their price is effectiveness:
without a label signal telling the model \emph{where} the trigger is, the model
must discover the trigger on its own, which requires either large triggers or
careful sample selection.

A recent line of work~\cite{generalcomponents2025} systematizes this with three
\emph{generalized components}---a quantization/blend trigger, a sample-selection
criterion, and a clean-label training protocol---and shows that the right
combination (e.g.\ a residual-squared selection) recovers competitive attack
success rates (ASR) on CIFAR-10. \textbf{However, the approach is brittle beyond
that comfort zone.} We reproduce these components faithfully and observe two
failure modes (Fig.~\ref{fig:main}): (1) on \emph{many-class / low-poison}
regimes (CIFAR-100 at $0.5\%$, Tiny-ImageNet at $0.25\%$) the strongest baseline
reaches only $\sim$44--85\% ASR; (2) on \emph{high-base-accuracy} tasks
(GTSRB, BA$\approx$100\%) the fixed trigger is diluted to $\le 34\%$ ASR---and
two of four baselines collapse to $0\%$. In all cases the benign task is so
easily learned that the tiny poison signal is washed out.

We argue the fix is to make the trigger \emph{adaptive and feature-aware} rather
than fixed. We propose \textbf{FAAT} (Feature-Aligned Adaptive Trigger), which
adds three components on top of the clean-label framework:
\begin{itemize}[leftmargin=*,topsep=2pt,itemsep=1pt]
  \item A \textbf{frozen global direction} $\dglobal$ (Narcissus-style~\cite{zeng2023narcissus})
        supplies a strong, transferable backdoor signal that is robust to benign dilution.
  \item A \textbf{bounded adaptive residual} $\dadapt$, optimized \emph{per poisoned
        image} in the DCT domain and clipped to a small $\ell_2$ budget, restores
        attack effectiveness without sacrificing imperceptibility.
  \item A \textbf{feature-alignment loss} $\Lalign$ pulls the representations of
        poisoned samples toward the target-class centroid, so that
        Activation-Clustering~\cite{chen2018activation} and
        Spectral-Signature~\cite{tran2018spectral} defenses cannot separate them
        from clean target-class samples.
\end{itemize}
All three are trained \emph{self-containedly}---FAAT does not depend on any
external pre-optimized noise, unlike Narcissus. Across CIFAR-10, CIFAR-100 and
Tiny-ImageNet, FAAT improves ASR by $+3.9$ to $+51.9$ points over the strongest
reproduced baseline (Tab.~\ref{tab:main}), preserves benign accuracy, stays
stealthy (SSIM $0.92$--$0.98$), and is not flagged by four sample-level defenses
(Tab.~\ref{tab:defense}). As a complementary contribution, \textbf{KST} is a
trigger variant whose DFT spectrum is \emph{structurally flat}
(peak/median $\approx\!1.0$ vs.\ $68.8$ for Narcissus), providing the first
clean-label trigger that defeats frequency-domain defenses at ASR parity
(Sec.~\ref{sec:kst}).

\paragraph{Contributions.} (1) FAAT, a self-contained clean-label backdoor with
a frozen global direction, a bounded DCT adaptive residual, and a
feature-alignment loss; (2) state-of-the-art clean-label ASR on three datasets,
especially on hard regimes where prior baselines collapse; (3) a systematic
defense evaluation showing FAAT evades sample-filtering defenses; (4) KST, a
flat-spectrum trigger that structurally evades frequency-domain
defenses.
